\chapter{Catégories triangulées.}
\VerdierPhysicalPageAnchor{0110}%
\VerdierSemanticAnchor{ch2}
\begingroup
\small
\setlength{\abovedisplayskip}{4pt plus 2pt minus 2pt}%
\setlength{\belowdisplayskip}{4pt plus 2pt minus 2pt}%

\VerdierSemanticAnchor{ch2-1}
\section{Définitions et premières propriétés.}

\VerdierSemanticAnchor{ch2-1-1}
\subsection{Définition des catégories triangulées.}

\VerdierSemanticAnchor{ch2-1-1-1}
\noindent\textbf{Définition 1.1.1.}\quad
Une \emph{catégorie triangulée} $\mathcal D$ est une $\mathbb Z$-catégorie
additive stricte (chap. I, 1.5.11) munie d'un ensemble de triangles
(chap. I, 3.2.1), appelés \emph{triangles distingués}, possédant les
propriétés suivantes :

\VerdierSemanticAnchor{ch2-1-1-1-tri}
\smallskip
\noindent\textbf{TRI :}
Tout triangle de $\mathcal D$ isomorphe à un triangle distingué est un
triangle distingué. Pour tout objet $X$ de $\mathcal D$, le triangle
\[
  X\xrightarrow{\operatorname{id}_X}X\longrightarrow 0
  \longrightarrow X[1]
\]
est distingué. Tout morphisme $u\colon X\longrightarrow Y$ de $\mathcal D$
est contenu dans un triangle distingué
\[
  X\xrightarrow{u}Y\xrightarrow{v}Z\xrightarrow{w}X[1].
\]

\VerdierSemanticAnchor{ch2-1-1-1-trii}
\noindent\textbf{TRII :}
Un triangle de $\mathcal D$
\[
  X\xrightarrow{u}Y\xrightarrow{v}Z\xrightarrow{w}X[1]
\]
est distingué si et seulement si le triangle
\[
  Y\xrightarrow{v}Z\xrightarrow{w}X[1]
  \xrightarrow{-u[1]}Y[1]
\]
est distingué.

\VerdierSemanticAnchor{ch2-1-1-1-triii}
\noindent\textbf{TRIII :}
Pour tout couple de triangles distingués :

\VerdierSemanticAnchor{ch2-1-1-1-triii-triangles}
\[
  \begin{aligned}
    X  &\xrightarrow{u}  Y  \xrightarrow{v}  Z  \xrightarrow{w}  X[1],\\
    X' &\xrightarrow{u'} Y' \xrightarrow{v'} Z' \xrightarrow{w'} X'[1]
  \end{aligned}
\]
et tout diagramme commutatif :

\VerdierSemanticAnchor{ch2-1-1-1-triii-square}
\begin{center}
  \begin{tikzcd}[column sep=4.5em,row sep=3em]
    X \arrow[r,"u"] \arrow[d,"f"'] & Y \arrow[d,"g"] \\
    X' \arrow[r,"u'"] & Y'\,,
  \end{tikzcd}
\end{center}
\endgroup
